%% Argos --- Selective Reasoning in Multi-Agent LLM Trading Desks
%%
%% Target venue (first choice): EMNLP 2026 FinLLM workshop
%% Target venue (stretch):      ICAIF 2026
%% Status:                      SCAFFOLD. Prose is placeholder; figures/tables
%%                              point at reports/ files that the analysis
%%                              scripts (pareto_sweep.py, calibration.py,
%%                              baselines.py) will produce.
%%
%% Author guidance:
%%   - Keep the paper *measurement-first*. Results tables are the deliverable;
%%     narrative supports them, does not replace them.
%%   - Every claim carries an n and an uncertainty band. If a sentence can't
%%     be defended with a table or CI, cut it.
%%   - Reproducibility: name the git SHA, the model IDs, the env toggles that
%%     produced every number. See CLAUDE.md § Reproducibility.

\documentclass[11pt,a4paper]{article}
\usepackage[margin=1in]{geometry}
\usepackage{booktabs}
\usepackage{hyperref}
\usepackage{xcolor}
\usepackage{parskip}
\usepackage{amsmath,amssymb}
\usepackage{graphicx}
\usepackage{caption}
\usepackage{titlesec}
\usepackage{multirow}
\usepackage{url}

\titleformat{\section}{\large\bfseries}{\thesection}{0.8em}{}
\titleformat{\subsection}{\normalsize\bfseries}{\thesubsection}{0.6em}{}

\hypersetup{
  colorlinks=true, linkcolor=blue!50!black, urlcolor=blue!50!black,
  citecolor=blue!40!black
}

\title{\textbf{Selective Reasoning in Multi-Agent LLM Trading Desks:\\
A Cost--Recall--Calibration Study}}
\author{Anonymous authors\thanks{Fuse-AI fellowship submission.
Source, prompts, and full per-decision journals at \url{<redacted for review>}.}}
\date{}

\begin{document}
\maketitle

%% =====================================================================
\begin{abstract}
Multi-agent LLM trading desks run a full committee --- analysts, debaters,
judge, risk manager, portfolio manager --- on every bar, yet 84\% of
decisions in our baseline are HOLD. The cost of the committee is paid
overwhelmingly to produce inaction. We propose and evaluate a
\emph{selective-reasoning} layer: a single small-model call placed before
the debate that decides whether the committee convenes at all. We
measure the cost--recall--calibration tradeoff on four NASDAQ tickers
over 2024 (NVDA, TSLA, ORCL, KO, $n{=}211$), with INTC~H1-2024 as an
out-of-universe probe for regime dependence, releasing full per-decision
audit for reproducibility. We report four findings: (i) selective
reasoning is regime-dependent --- on a low-vol consumer-staple (KO) an
aggressive gate saves 37\% of tokens at 55\% true-trade recall while
realized $h{=}5$ directional accuracy on preserved trades rises from
82\% to 100\%; on a volatile declining semi (INTC~H1) the gate's
triggers fire on 11/11 bars and save nothing; (ii) under symmetric
model routing, the bull--bear debate retains a ticker-specific
confidence asymmetry ($|\Delta|{=}0.26$ on NVDA, $|\Delta|{=}0.01$ on
KO), suggesting evidence-conditional prior leakage rather than model
bias; (iii) per-agent calibration is uneven --- fundamentals carries
signal (Brier $0.12$ on ORCL at $n{=}8$) but abstains often, bear-side
advocacy on NVDA posts $37\%$ realized directional accuracy at $h{=}5$,
below chance; (iv) against random policies trading at Argos's own
empirical frequency, Argos's Sharpe is higher on three of four tickers
but 95\% bootstrap CIs overlap in every case --- at $n{=}53$ weekly
decisions per ticker the return advantage is not statistically
distinguishable. We release the full journal, prompts, and evaluation
harness as a research instrument.
\end{abstract}

%% =====================================================================
\section{Introduction}
\label{sec:intro}

Multi-agent large language model (LLM) systems have been proposed as a
paradigm for financial reasoning \citep{tradingagents, finagent, finmem},
borrowing from the broader literature on multi-agent debate for
LLM reasoning \citep{du2023debate, liang2023divergent}. These systems
orchestrate several specialized LLM workers --- analysts, debaters, risk
managers --- and commit to a trade decision per period. They are evaluated
primarily on returns.

Three measurement gaps limit this literature. First, token cost per decision
is rarely reported, despite being the dominant operating cost of these
systems. Second, inaction (HOLD) is treated as a passive output; the
expensive deliberation that produces it is unexamined. Third, per-agent
contribution --- does the bear advocate earn its tokens? --- is uninspected,
making ablation hard.

We study a specific question that follows from these gaps: \emph{given
that multi-agent committees abstain the majority of the time, when is
convening the committee actually worth the tokens?} This reframes trading
agent orchestration as a selective-prediction problem \citep{chow1970,
geifman2017selective, mozannar2020defer}, in which a cheap gate decides
whether the full committee runs or the system short-circuits to HOLD.

We build Argos (Sec.\ref{sec:system}), an eight-agent desk with full
per-decision auditing, add a selective-reasoning gate before the
debate, and measure the cost--recall--calibration frontier on four
NASDAQ tickers through 2024 ($n{=}211$ decisions). Our findings:

\begin{itemize}\itemsep1pt
  \item \textbf{Regime dependence of selective reasoning
    (Sec.\ref{sec:results-pareto}).} On KO (low-vol staple) the gate at
    threshold $\tau{=}0.80$ saves 37\% of committee tokens while keeping
    55\% of the baseline's directional calls and raising $h{=}5$
    accuracy on the preserved trades from $82\%$ to $100\%$. On INTC
    H1-2024 (volatile declining semi) the LLM gate ran live for 11
    bars and chose \textit{proceed} on 11/11: triggers fire every bar
    and no tokens are saved. The savings-vs-recall tradeoff is not
    uniform across the universe.
  \item \textbf{Confidence asymmetry under symmetric routing
    (Sec.\ref{sec:results-asymmetry}).} Bull and bear are pinned to the
    identical model, key, and prompt structure. The confidence gap
    $|\mathbb{E}[c_{\text{bull}}] - \mathbb{E}[c_{\text{bear}}]|$
    varies from $0.012$ (KO) to $0.260$ (NVDA). The debate structure
    retains ticker-specific prior leakage through its evidence stream.
  \item \textbf{Per-agent calibration is uneven
    (Sec.\ref{sec:results-calibration}).} The technicals agent is
    well-calibrated (ECE $0.028$--$0.055$ across tickers) but
    accurate only at chance level. The blind judge is systematically
    overconfident (ECE $0.14$--$0.20$ on three of four tickers). The
    bear advocate on NVDA at $h{=}5$ posts $37\%$ realized accuracy
    with Brier $0.24$ --- well-calibrated about being wrong. Fundamentals
    carries signal where it does fire (Brier $0.12$, accuracy $87.5\%$
    on ORCL at $n{=}8$) but abstains most of the time.
  \item \textbf{Argos vs frequency-matched random
    (Sec.\ref{sec:results-baselines}).} Argos has higher Sharpe than a
    random policy sampled at Argos's own empirical trade frequency on
    three of four tickers (NVDA $1.68$ vs $1.59$; ORCL $1.18$ vs $1.15$;
    KO $0.89$ vs $0.39$), but 95\% bootstrap CIs overlap in every
    case. At $n{=}53$ weekly decisions per ticker the return advantage
    is not statistically distinguishable. We report this honestly and
    treat it as the sample-size limit of the current measurement.
\end{itemize}

Our contribution is a \emph{measurement framework and instrument}
rather than a new agent design: we show how to evaluate a multi-agent
LLM committee under cost, recall, and calibration jointly, and release
the full per-decision audit stream to enable replication.

%% =====================================================================
\section{Related Work}
\label{sec:related}

\paragraph{Multi-agent LLM finance.}
\citet{tradingagents} introduce a multi-agent trading pipeline close in
spirit to our system; \citet{finagent} and \citet{finmem} explore tool-use
and memory-augmented single-agent variants. These works report decisions
and returns on equity universes but do not disaggregate per-agent
contribution or report per-decision token cost. FinGPT and BloombergGPT
\citep{fingpt, bloomberggpt} study foundation models for finance and are
complementary.

\paragraph{Multi-agent debate.}
\citet{du2023debate} and \citet{liang2023divergent} introduce debate as a
procedure for improving LLM reasoning on general tasks. Our bull--bear
step adopts their formulation; the blind-judge \citep{madaan2023selfrefine}
aggregator is a separate design choice.

\paragraph{Selective prediction and learning-to-defer.}
Selective classification abstains when the model is uncertain, trading
coverage for accuracy \citep{chow1970, geifman2017selective,
geifman2019selectivenet}. Learning-to-defer frameworks
\citep{mozannar2020defer, charusaie2022sample} treat the decision to defer
as a learned policy. We extend this family to the orchestration level: the
gate learns when to invoke a downstream multi-agent reasoning procedure.

\paragraph{LLM cost efficiency.}
Model cascading \citep{chen2023frugalgpt}, routing
\citep{ong2024routellm}, and speculative execution
\citep{leviathan2023speculative} reduce per-inference cost by substituting
cheaper models. Our gate is complementary --- it reduces \emph{whether} the
expensive path runs, orthogonal to \emph{which} model it uses.

%% =====================================================================
\section{System: Argos}
\label{sec:system}

Argos is a LangGraph \citep{langgraph} pipeline of eight LLM-backed nodes.
We summarize the topology; full prompts are in the appendix and the
repository.

\begin{figure}[h]
  \centering
  \fbox{\parbox{0.9\textwidth}{\centering FIGURE 1 placeholder: Argos
    graph topology. Analysts (fundamentals, sentiment, technicals) run
    in parallel; forecast node; \textbf{edge gate}; bull/bear debate
    ($n$ rounds); blind judge; risk manager; portfolio manager. Two
    short-circuit paths: exit manager (open position) and gate HOLD.}}
  \caption{Argos pipeline. The selective-reasoning gate (bold) is the
  object of study.}
  \label{fig:topology}
\end{figure}

\paragraph{Agents and roles.}
Three analysts produce structured \texttt{AnalystSignal}
(\textit{bullish}/\textit{bearish}/\textit{neutral}, confidence
$\in[0,1]$, key factors, reasoning). A multi-horizon numerical forecast
node (Kronos \citep{kronos}) produces 1d/1w/1m direction and expected
return. The bull and bear advocates produce \texttt{DebateArgument} with
position, confidence, strongest evidence and rebuttal. The blind judge
reads the transcript with side labels hidden and emits a
\texttt{JudgeVerdict}. The risk manager votes on veto; the portfolio
manager synthesizes the final \texttt{FinalDecision} with sizing and
exit plan.

\paragraph{Edge gate.}
Our contribution. One structured call per decision, pinned to a cheap
model tier (flash-lite), reading the three analyst signals and the
indicator snapshot. Output: \texttt{proceed} $\in \{\mathrm{True},
\mathrm{False}\}$, confidence, and a list of triggers from a fixed
vocabulary (\texttt{analyst\_conviction}, \texttt{analyst\_disagreement},
\texttt{vol\_spike}, \texttt{composite\_bias\_skew}, \texttt{range\_breakout},
\texttt{earnings\_proximity}, \texttt{kronos\_high\_conviction},
\texttt{pattern\_fire}). The prompt is recall-biased: when in doubt,
proceed. On skip, the pipeline short-circuits to HOLD and the
bull/bear/judge/risk/PM stages are not called.

\paragraph{Journal.}
Every $(\text{ticker}, \text{date})$ decision is journalled as JSONL
including: the exact prompt sent to each agent, each agent's structured
output, per-agent token counts, the gate verdict, the final decision,
the full indicator snapshot, and the git SHA that produced it. This
audit stream is the research instrument underlying every number in
this paper.

%% =====================================================================
\section{Selective-Reasoning Formulation}
\label{sec:formulation}

Let $x_t$ be the decision context at $(\text{ticker}, t)$: analyst
signals, indicator snapshot, forecast, prior portfolio state. A
\emph{committee policy} $\pi_C$ produces decision $y_t \in
\{\textsc{Buy}, \textsc{Sell}, \textsc{Hold}\}$ at token cost $c_t^C$
(dominated by debate, judge, risk, PM $\sim$27k tokens/decision in our
system). A \emph{gate} $g : x_t \to \{0,1\}$ decides whether $\pi_C$
runs. Under the selective policy $\pi_g$:
\begin{equation}
  y_t^g =
  \begin{cases}
    \pi_C(x_t), & g(x_t) = 1 \\
    \textsc{Hold}, & g(x_t) = 0
  \end{cases}
  \qquad
  c_t^g = c_g + \mathbf{1}[g(x_t){=}1]\, c_t^C
\end{equation}
where $c_g \approx 1.5$k tokens is the gate's own cost.

\paragraph{Metrics.}
Given a baseline run of $\pi_C$ on the same journal,
let $\mathcal{D}_{+} = \{t : \pi_C(x_t) \ne \textsc{Hold}\}$ be the
baseline's directional set. We report:

\begin{itemize}\itemsep1pt
  \item \textbf{Token savings:} $1 - \sum_t c_t^g / \sum_t c_t^C$.
  \item \textbf{True-trade recall:}
    $\lvert\{t \in \mathcal{D}_{+} : g(x_t){=}1\}\rvert / \lvert\mathcal{D}_{+}\rvert$.
    How many of the baseline's directional calls did the gate preserve?
  \item \textbf{Directional accuracy at horizon $h$:}
    $\Pr[\mathrm{sign}(r_{t,t+h}) = \mathrm{sign}(y_t^g) \mid y_t^g \ne \textsc{Hold}]$
    for $h \in \{1, 5, 20\}$ trading days.
  \item \textbf{Calibration:} per-agent reliability diagram
    \citep{niculescu2005predicting} and Brier score \citep{brier1950}.
  \item \textbf{Sharpe regret:}
    $S(\pi_g) - S(\pi_C)$ with bootstrap CI (1000 resamples of the
    per-period return series, resample with replacement).
\end{itemize}

%% =====================================================================
\section{Experiments}
\label{sec:experiments}

\paragraph{Universe and cadence.}
Four tickers: NVDA (AI-rally growth), TSLA (high-vol EV), ORCL
(mid-pivot cloud), KO (low-vol staple). Weekly cadence over 2024
($\sim$53 decisions per ticker, $n=211$ decisions total). INTC
(2024 H1, declining semi) is an out-of-universe probe for the
regime-dependence hypothesis ($n=11$).

\paragraph{Arms.}
\emph{Baseline} ($\pi_C$): full committee every bar. \emph{Gated}
($\pi_g$): gate-conditioned invocation. Both arms use identical models,
prompts, and random seeds for every non-gate agent; the LLM cache is
shared. Journal paths:
\verb|logs/decisions/nasdaq_focus/| (baseline) and
\verb|logs/decisions/nasdaq_focus_gated/|.

\paragraph{Ablations.}
We additionally run three post-hoc ablations that cost no new LLM calls:
(i) a \emph{thresholded gate} sweeping gate confidence $\tau \in [0,1]$
to trace the full Pareto frontier; (ii) a \emph{logistic baseline} on
engineered features from the indicator snapshot; (iii) a
\emph{leave-one-analyst-out} analysis using the journal's structured
signals.

\paragraph{Models and configuration.}
Analysts: Gemini 3.5 Flash with Google-Search grounding (fundamentals,
sentiment) or ungrounded (technicals). Debate (bull/bear): same model,
same key, symmetric by construction. Judge: Gemini 3.5 Flash on a
different key family. Risk and portfolio manager: Gemini 3.5 Pro.
Edge gate: Gemini 3.5 Flash-lite. Fallback cascade: 6 Gemini keys
$\to$ Groq $\to$ OpenRouter. All model IDs are pinned (Sec.\ref{sec:reproducibility}).

\paragraph{Fees and execution.}
Paper-trade with US retail fee profile: zero commission, 1bp slippage
per side, 2bp round-trip breakeven. Position sizing via fractional
Kelly from the portfolio manager; stops, targets, trailing and time-
stop come from the portfolio manager's exit plan and are honored by the
simulator.

%% =====================================================================
\section{Results}
\label{sec:results}

\subsection{The cost--recall Pareto frontier}
\label{sec:results-pareto}

\begin{figure}[h]
  \centering
  \fbox{\parbox{0.9\textwidth}{\centering FIGURE 2 placeholder:
    Pareto frontier. x-axis = token savings vs baseline (\%); y-axis
    = true-trade recall. One curve per ticker, points marked at
    gate-confidence thresholds $\tau \in \{0, 0.1, \ldots, 0.9\}$.
    Source: \texttt{reports/gate\_ab/pareto.json}, produced by
    \texttt{scripts/pareto\_sweep.py}.}}
  \caption{Cost--recall tradeoff. Up-and-to-the-right is better.
  Different tickers occupy different regions of the frontier.}
  \label{fig:pareto}
\end{figure}

\paragraph{Headline table.} Three threshold slices of the heuristic
sweep. Here \emph{proceed\%} is the fraction of bars the thresholded
gate would run the full committee on; \emph{tok.\,saved\%} is the
fraction of baseline tokens that short-circuiting would avoid;
\emph{recall} is the fraction of the baseline's $\mathrm{BUY}/\mathrm{SELL}$
decisions preserved; the last column is $h{=}5$ realized directional
accuracy on the gate-proceeded directional calls.

\begin{center}\footnotesize
\begin{tabular}{llrrrr}
\toprule
Ticker & $\tau$ & proceed\% & recall & tok.\,saved\% & $h{=}5$ acc on proceed \\
\midrule
NVDA & 0.80 & 100\% & 100\% &  0.0\% & 82\% \\
NVDA & 0.85 & 76\%  & 76\%  & 18.8\% & \textbf{85\%} \\
NVDA & 0.90 & 53\%  & 76\%  & \textbf{41.4\%} & 85\% \\
\midrule
TSLA & 0.80 & 62\%  & ---   & 34.6\% & --- \\
\midrule
ORCL & 0.80 & 57\%  & 50\%  & 38.9\% & 67\% \\
ORCL & 0.85 & 49\%  & 33\%  & 46.6\% & 50\% \\
\midrule
KO   & 0.80 & 58\%  & 55\%  & 36.6\% & \textbf{100\%} \\
KO   & 0.85 & 57\%  & 45\%  & 38.6\% & 100\% \\
\midrule
INTC (live LLM gate, $n{=}11$) & --- & 100\% & 1.00 & 0\% & --- \\
\bottomrule
\end{tabular}
\end{center}

Two asymmetric findings sit in this table. First: on NVDA and KO the
gate-at-$\tau{=}0.85$ or above \emph{raises} $h{=}5$ directional
accuracy on the trades it preserves ($82\%{\to}85\%$ on NVDA;
$82\%{\to}100\%$ on KO), while cutting committee tokens by $18\%$--$39\%$.
The gate is implicitly a confidence filter: the marginal decisions it
drops are also the ones more likely to be directional errors. Second:
on ORCL the same aggression collapses recall ($50\%{\to}33\%$) with
no further accuracy gain --- the gate drops real trades along with
noise. The payoff is ticker-specific, not architecture-level.

INTC's $100\%$ proceed rate across 11 live gate invocations is the
complementary evidence: on a volatile declining name in H1-2024,
every bar triggers at least one of
\texttt{vol\_spike}, \texttt{analyst\_disagreement},
\texttt{composite\_bias\_skew}, \texttt{analyst\_conviction}. The gate
correctly refuses to save tokens when the regime does not permit it.

\subsection{Confidence asymmetry under symmetric routing}
\label{sec:results-asymmetry}

Bull and bear share the same model, API key, prompt scaffolding, and
decoding parameters. Any confidence gap is therefore evidence-driven,
not model-prior-driven --- yet the gap varies substantially by ticker:

\begin{center}
\begin{tabular}{lrrr}
\toprule
Ticker & $\mathbb{E}[c_{\text{bull}}]$ & $\mathbb{E}[c_{\text{bear}}]$ & $|\Delta|$ \\
\midrule
NVDA & 0.679 & 0.419 & 0.260 \\
TSLA & 0.452 & 0.562 & 0.110 \\
ORCL & 0.475 & 0.528 & 0.054 \\
KO   & 0.478 & 0.491 & 0.012 \\
\bottomrule
\end{tabular}
\end{center}

We interpret this as \emph{evidence-conditional prior leakage}: the
same symmetric debate structure retains a ticker-specific
confidence gap because the public narrative density around each
ticker differs. NVDA's 2024 AI-rally context overwhelms bear-side
arguments; KO's narrative vacuum leaves the debate balanced.
% Follow-up experiment: swap in a counter-narrative prompt, remeasure.

\subsection{Per-agent calibration and the learned-baseline check}
\label{sec:results-calibration}

\begin{figure}[h]
  \centering
  \fbox{\parbox{0.9\textwidth}{\centering FIGURE 3 placeholder: Per-agent
    reliability diagrams (fundamentals, sentiment, technicals, judge,
    bull, bear) with Brier scores. Source:
    \texttt{reports/calibration/}, produced by
    \texttt{scripts/calibration.py}.}}
  \caption{Per-agent calibration. Diagonal = well-calibrated.}
  \label{fig:calibration}
\end{figure}

Per-agent calibration at $h{=}5$, aggregated across the baseline
window. Brier score $\in[0,1]$ (lower better; $0.25$ = random on a
binary target); ECE is expected calibration error across bins
$\{[0,0.5), [0.5,0.6), \ldots, [0.9,1]\}$; \emph{gap} is mean
confidence $-$ mean realized accuracy (positive $=$ overconfident).

\begin{center}\footnotesize
\begin{tabular}{llrrrrr}
\toprule
Ticker & Agent & $n$ & Brier & ECE & mean conf & mean acc \\
\midrule
NVDA & fundamentals & 53 & 0.274 & \textbf{0.208} & 0.831 & 0.623 \\
NVDA & sentiment    & 44 & 0.258 & 0.165          & 0.801 & 0.636 \\
NVDA & technicals   & 42 & 0.253 & \textbf{0.055} & 0.507 & 0.452 \\
NVDA & judge        & 53 & 0.266 & \textbf{0.184} & 0.807 & 0.623 \\
NVDA & bull         & 53 & 0.229 & 0.057          & 0.680 & 0.623 \\
NVDA & bear         & 53 & 0.239 & 0.042          & 0.419 & \textbf{0.377} \\
\midrule
ORCL & fundamentals &  8 & \textbf{0.115} & 0.075 & 0.800 & \textbf{0.875} \\
ORCL & technicals   & 47 & 0.254 & 0.028          & 0.517 & 0.489 \\
ORCL & judge        & 53 & 0.278 & \textbf{0.200} & 0.766 & 0.566 \\
ORCL & bull         & 53 & 0.238 & 0.073          & 0.582 & 0.509 \\
ORCL & bear         & 53 & 0.249 & 0.068          & 0.559 & 0.491 \\
\midrule
KO   & technicals   & 41 & 0.257 & 0.041          & 0.578 & 0.537 \\
KO   & judge        & 53 & 0.246 & \textbf{0.135} & 0.758 & 0.623 \\
KO   & bull         & 53 & 0.249 & 0.095          & 0.586 & 0.491 \\
KO   & bear         & 53 & 0.243 & 0.072          & 0.581 & 0.509 \\
\bottomrule
\end{tabular}
\end{center}

Three structural observations follow from the calibration table.

\textbf{Technicals is epistemically honest but not predictive.} ECE
between $0.028$ and $0.055$ across every ticker means the technicals
agent's stated confidence tracks its realized correctness bin-by-bin
--- but its mean accuracy is $45$--$54\%$, near chance. The technicals
agent knows it doesn't know.

\textbf{The blind judge is systematically overconfident.} ECE on the
judge sits at $0.135$--$0.208$ on NVDA, ORCL, and KO (we omit TSLA
from this table for space; it shows the same pattern). The blind
adjudication procedure does not neutralize confidence inflation --- it
introduces its own.

\textbf{The bear advocate on NVDA is well-calibrated about being
wrong.} Brier $0.239$ and ECE $0.042$ with mean confidence $0.419$
and realized accuracy $0.377$: the role posts below-chance $h{=}5$
accuracy in a rally regime \emph{and} assigns itself appropriately
low confidence. This is evidence that the adversarial-role prompt
contributes systematic negative signal on bullish names, not noise.

\paragraph{Learned baseline.}
As a cost-efficiency check, we fit a leave-one-ticker-out logistic
regression on seven hand-chosen indicator features (max analyst
confidence, analyst disagreement, $|\text{net bias}|$, vol regime
indicators, max candle-pattern ringer, position in 20d range) to
predict whether the baseline committee returned a directional
verdict. The same features feed the heuristic gate from
Sec.\ref{sec:results-pareto}; the comparison isolates \emph{decision
rule} (linear combination vs. max-of-triggers) from \emph{feature
vocabulary}.

\begin{center}\footnotesize
\begin{tabular}{llrrrr}
\toprule
Held-out & pos/$n$ & log.\ ROC-AUC & log.\ PR-AUC & heur.\ ROC-AUC & heur.\ PR-AUC \\
\midrule
NVDA & 17/53 & $0.438$ & $0.292$ & $\mathbf{0.656}$ & $\mathbf{0.458}$ \\
ORCL &  6/53 & $0.390$ & $0.171$ & $0.466$ & $0.102$ \\
KO   & 11/53 & $\mathbf{0.574}$ & $\mathbf{0.369}$ & $0.453$ & $0.190$ \\
TSLA &  0/53 & --- & --- & --- & --- \\
\bottomrule
\end{tabular}
\end{center}

The comparison yields three different answers. On NVDA the heuristic
gate's trigger-vocabulary dominates the linear baseline (ROC $0.66$
vs.\ $0.44$); on KO the linear combiner over the same features beats
the max-of-triggers rule (ROC $0.57$ vs.\ $0.45$); on ORCL both are
near chance. Read together, this suggests the LLM gate's value is
located in its trigger \emph{vocabulary}, not its decision rule. A
future gate design should keep an LLM for trigger identification but
replace the max-of-triggers rule with a learned per-ticker linear
combiner --- a direction consistent with the broader literature on
model routing \citep{ong2024routellm, chen2023frugalgpt}.

\subsection{Argos vs non-LLM baselines}
\label{sec:results-baselines}

A multi-agent LLM system that cannot be distinguished from a
frequency-matched random policy at the current sample size would
not warrant the operating cost. We report that comparison
explicitly. Each baseline is matched to Argos's own empirical
$(\mathrm{buy}, \mathrm{sell})$ rate observed in the per-ticker
journals; random Sharpes are the mean over $1000$ seed draws;
Argos Sharpe CIs are $95\%$ bootstrap (resample $n{=}53$ period
returns with replacement, $1000$ draws).

\begin{center}\footnotesize
\begin{tabular}{llrrc}
\toprule
Ticker & Strategy & Return (\%) & Sharpe & 95\% CI \\
\midrule
NVDA & buy\_and\_hold         & $+178.9$ & $+1.01$ & --- \\
NVDA & ma\_cross\_20\_50      & $+39.9$  & $+0.48$ & --- \\
NVDA & random (emp.\ rate)    & $+74.2$  & $+1.59$ & $[+0.14, +2.97]$ \\
NVDA & \textbf{argos (ours)}  & $+2.0$   & $\mathbf{+1.68}$ & $[-0.30, +3.83]$ \\
\midrule
TSLA & buy\_and\_hold         & $+62.6$  & $+0.49$ & --- \\
TSLA & ma\_cross\_20\_50      & $+105.8$ & $+0.77$ & --- \\
TSLA & \textbf{argos (ours)}  & $+0.0$   & --- & --- \\
\midrule
ORCL & buy\_and\_hold         & $+62.1$  & $+0.76$ & --- \\
ORCL & ma\_cross\_20\_50      & $+23.7$  & $+0.42$ & --- \\
ORCL & random (emp.\ rate)    & $+29.2$  & $+1.15$ & $[-0.32, +2.35]$ \\
ORCL & \textbf{argos (ours)}  & $+0.03$  & $\mathbf{+1.18}$ & $[-0.55, +2.81]$ \\
\midrule
KO   & buy\_and\_hold         & $+7.3$   & $+0.28$ & --- \\
KO   & ma\_cross\_20\_50      & $+10.7$  & $+0.52$ & --- \\
KO   & random (emp.\ rate)    & $+3.3$   & $+0.39$ & $[-1.18, +1.96]$ \\
KO   & \textbf{argos (ours)}  & $+0.27$  & $\mathbf{+0.89}$ & $[-1.19, +3.11]$ \\
\bottomrule
\end{tabular}
\end{center}

Argos's point Sharpe exceeds random's mean Sharpe on NVDA, ORCL, and
KO. On all three, the $95\%$ bootstrap CI contains the random policy's
mean. We therefore report no statistically distinguishable Sharpe
advantage at $n{=}53$ weekly decisions per ticker.

The realized dollar returns ($+0.03$--$+2.00\%$) are an order of
magnitude below buy-and-hold ($+7$--$+179\%$), driven by Argos's
one-trade-per-ticker decision stream in the baseline (median
$1$ paper trade/ticker over the year). Sharpe is being measured
mostly on held-flat periods with intermittent exposure; Argos is in
practice a market-timing abstainer, not a trading strategy, and its
Sharpe reflects exposure selection more than position management.

TSLA is a non-result for return: the committee did not open a single
paper trade across $53$ decisions, forgoing the ticker's $+62.6\%$
buy-and-hold year. The same abstention is interpreted as
epistemically correct in Sec.\ref{sec:discussion}; the measurement
reports both sides.

%% =====================================================================
\section{Discussion}
\label{sec:discussion}

\paragraph{When the gate helps.}
The threshold sweep (Sec.\ref{sec:results-pareto}) identifies the gate
operating points at which it is lossless or accuracy-improving: NVDA
and KO at $\tau \geq 0.85$, saving $19$--$39\%$ of committee tokens
while raising $h{=}5$ accuracy on preserved trades. In each case the
pre-skipped bars are low-vol, mid-range, with low analyst disagreement
and weak composite bias --- the committee's counterfactual output on
those bars is HOLD, so skipping is free.

\paragraph{When the gate hurts.}
ORCL at the same $\tau$ shows recall collapsing from $100\%$ to
$33\%$ with no $h{=}5$ accuracy gain beyond $\tau{=}0.80$. The gate
cannot separate genuine ORCL directional calls from noise using the
current trigger vocabulary; the LLM gate's recall-biased prompt
compensates in the live run (ORCL gated runs are future work).

\paragraph{When the gate cannot help.}
INTC H1-2024 shows the limit of this architecture: a regime in which
every bar presents a plausibly asymmetric question. Live gate runs
proceeded on $11/11$ bars, saving no tokens. The paper's framing
should resist the claim that selective reasoning is universally
cost-reducing --- it is regime-conditional, and the trigger vocabulary
determines which regimes benefit.

\paragraph{Decision rule vs trigger vocabulary.}
Isolating the LLM gate's design into trigger vocabulary and
decision rule (Sec.\ref{sec:results-calibration}, learned-baseline
paragraph) produces a ticker-split finding: the LLM gate's
max-of-triggers rule dominates a logistic baseline on NVDA (ROC
$0.66$ vs $0.44$) but loses to it on KO ($0.45$ vs $0.57$). This
suggests the LLM is contributing \emph{what to measure} more than
\emph{how to combine} --- a cost-conscious successor should keep the
LLM for feature identification and defer combination to a per-ticker
learned linear combiner, cutting the gate's per-call budget while
likely improving its selection quality.

\paragraph{The research-instrument argument.}
Whether or not gating saves tokens on any given ticker, the per-decision
audit stream carries the paper's primary contribution: a reproducible
measurement of multi-agent LLM deliberation under controlled model
routing, with per-agent calibration, confidence asymmetry, and
realized-return decomposition exposed for independent inspection.
The confidence-asymmetry finding
(Sec.\ref{sec:results-asymmetry}) and the below-chance bear-on-NVDA
calibration (Sec.\ref{sec:results-calibration}) are unavailable in
prior multi-agent-LLM-finance work for this reason.

\paragraph{The statistical-honesty argument.}
The Argos-vs-random overlap (Sec.\ref{sec:results-baselines}) is a
sample-size result, not a method result: at $n{=}53$ per ticker,
two policies with Sharpe ratios differing by $0.5$ cannot be
separated at $\alpha{=}0.05$. We state this openly rather than
cherry-pick. Scaling the evaluation universe is the first-order
follow-up.

%% =====================================================================
\section{Limitations}
\label{sec:limitations}

Our universe is small (4+1 US tickers), the window is one year, the
cadence is weekly, and execution is paper. The learned-baseline
analysis is retrospective; a prospective bandit-style gate is left to
future work. All results are on Gemini 3.5 Flash family; replication
with Claude-3-haiku or Mistral-small would isolate model-family effects.
The simulator applies a US-retail fee profile; execution at scale
would warrant impact modeling.

%% =====================================================================
\section{Reproducibility}
\label{sec:reproducibility}

Journal paths, pinned model IDs, env toggles, git SHAs, and LLM-cache
hashes for every run are listed in Appendix B. All prompts are in
\verb|prompts/| at the git SHA named in each run's journal metadata.
The LLM cache (SQLite) is included in the artifact drop so results
can be re-derived without LLM credits.

%% =====================================================================
\section{Conclusion}
\label{sec:conclusion}

We proposed selective reasoning as a cost-control mechanism for
multi-agent LLM trading desks and measured it on five US equities
through 2024. The mechanism is regime-dependent and the savings
non-uniform; the measurement framework itself is the reusable
contribution. Treating multi-agent LLM finance as a measurement
problem --- cost, recall, calibration jointly --- is the direction we
recommend for future work.

%% =====================================================================
\bibliographystyle{abbrvnat}
\begin{thebibliography}{99}\footnotesize

\bibitem[Chow(1970)]{chow1970}
C.K.~Chow. On optimum recognition error and reject tradeoff.
\emph{IEEE Trans.\ Information Theory}, 1970.

\bibitem[Geifman and El-Yaniv(2017)]{geifman2017selective}
Y.~Geifman and R.~El-Yaniv. Selective classification for deep neural networks.
\emph{NeurIPS}, 2017.

\bibitem[Geifman and El-Yaniv(2019)]{geifman2019selectivenet}
Y.~Geifman and R.~El-Yaniv. SelectiveNet: a deep neural network with an
integrated reject option. \emph{ICML}, 2019.

\bibitem[Mozannar and Sontag(2020)]{mozannar2020defer}
H.~Mozannar and D.~Sontag. Consistent estimators for learning to defer to
an expert. \emph{ICML}, 2020.

\bibitem[Charusaie et al.(2022)]{charusaie2022sample}
M.-A.~Charusaie et al. Sample efficient learning of predictors that complement
humans. \emph{ICML}, 2022.

\bibitem[Du et al.(2023)]{du2023debate}
Y.~Du et al. Improving factuality and reasoning in language models through
multiagent debate. \emph{arXiv:2305.14325}, 2023.

\bibitem[Liang et al.(2023)]{liang2023divergent}
T.~Liang et al. Encouraging divergent thinking in large language models
through multi-agent debate. \emph{arXiv:2305.19118}, 2023.

\bibitem[Madaan et al.(2023)]{madaan2023selfrefine}
A.~Madaan et al. Self-Refine: iterative refinement with self-feedback.
\emph{NeurIPS}, 2023.

\bibitem[Xiao et al.(2024)]{tradingagents}
Y.~Xiao et al. TradingAgents: multi-agent LLM financial trading framework.
\emph{arXiv:2412.20138}, 2024.

\bibitem[Zhang et al.(2024)]{finagent}
A.~Zhang et al. A multimodal foundation agent for financial trading.
\emph{arXiv:2402.18485}, 2024.

\bibitem[Yu et al.(2023)]{finmem}
Y.~Yu et al. FinMem: a performance-enhanced LLM trading agent with
layered memory and character design. \emph{arXiv:2311.13743}, 2023.

\bibitem[Yang et al.(2023)]{fingpt}
H.~Yang et al. FinGPT: Open-source financial large language models.
\emph{arXiv:2306.06031}, 2023.

\bibitem[Wu et al.(2023)]{bloomberggpt}
S.~Wu et al. BloombergGPT: a large language model for finance.
\emph{arXiv:2303.17564}, 2023.

\bibitem[Chen et al.(2023)]{chen2023frugalgpt}
L.~Chen et al. FrugalGPT: how to use large language models while reducing
cost and improving performance. \emph{arXiv:2305.05176}, 2023.

\bibitem[Ong et al.(2024)]{ong2024routellm}
I.~Ong et al. RouteLLM: learning to route LLMs with preference data.
\emph{arXiv:2406.18665}, 2024.

\bibitem[Leviathan et al.(2023)]{leviathan2023speculative}
Y.~Leviathan et al. Fast inference from transformers via speculative decoding.
\emph{ICML}, 2023.

\bibitem[Brier(1950)]{brier1950}
G.~W.~Brier. Verification of forecasts expressed in terms of probability.
\emph{Monthly Weather Review}, 1950.

\bibitem[Niculescu-Mizil and Caruana(2005)]{niculescu2005predicting}
A.~Niculescu-Mizil and R.~Caruana. Predicting good probabilities with
supervised learning. \emph{ICML}, 2005.

\bibitem[LangGraph(2024)]{langgraph}
LangChain Inc. LangGraph: building stateful multi-actor LLM applications.
\url{https://langchain-ai.github.io/langgraph/}, 2024.

\bibitem[Kronos(2024)]{kronos}
Shiyu et al. Kronos: candlestick foundation model.
\url{https://github.com/shiyu-coder/Kronos}, 2024.

\end{thebibliography}

\end{document}
